\documentclass[ ../main.tex]{subfiles}
\providecommand{\mainx}{..}
\begin{document}
\section{Introduction}

The \emph{Bernoulli set}~\cite{bsets} is a probabilistic model of a set with
quantifiable false positive and false negative rates. A \emph{positive}
Bernoulli set has a false negative rate of zero: every true member is
recognized, but non-members may be incorrectly accepted at a specified false
positive rate~$\varepsilon$. This model captures the semantics of data
structures such as the Bloom filter~\cite{bf}, which implements the positive
Bernoulli set abstract data type.

In this paper, we study the \emph{Perfect Hash Filter} (\gls{gls-phf}): a data
structure that implements the static positive Bernoulli set over $\BitSet^*$ by
placing each member at a unique slot, via a minimal perfect hash function, and
storing a short fingerprint of the member at that slot. Membership is tested by
fingerprint equality. This minimal-perfect-hash-plus-fingerprint structure is
the well-known static approximate-membership-query construction: it is the
mechanism behind the Bloomier filter~\cite{bloomier}, and close relatives appear
in xor~\cite{graf2020}, ribbon~\cite{dillinger2021}, and cuckoo~\cite{fan2014}
filters. The base structure is therefore not new. Our contribution is its
treatment within the Bernoulli framework: realizing it as a Bernoulli set ADT,
extending it to a two-rate structure, lifting it to a key-value map, and tying
its measured space to the framework's space-accuracy bound. \Cref{sec:related}
situates the \gls{gls-phf} against this prior work in detail.

The \gls{gls-phf} is \emph{succinct}: its expected space complexity approaches
the information-theoretic lower bound of $-\log_2 \varepsilon$~bits per element
established by Carter et al.~\cite{carter1978} (\cref{thm:approx_l_b}),
asymptotically as $\varepsilon \to 0$. It does not attain the bound at finite
$\varepsilon$. By comparison, the Bloom filter spends $-1.44 \log_2 \varepsilon$
bits per element, a factor of $1/\ln 2 \approx 1.44$ above the bound. For false
positive rates $\varepsilon < 2^{-3}$, the idealized \gls{gls-phf} is the more
space-efficient of the two.

\paragraph{Contributions.}
\begin{enumerate}
    \item We realize the \gls{gls-phf}, the minimal-perfect-hash-plus-fingerprint
    structure, as a faithful implementation of the \emph{positive Bernoulli
    set}~\cite{bsets} (false positive rate $\varepsilon = 2^{-M}$, zero false
    negatives), so that it inherits the family's error model, composition theory,
    and space bounds automatically (\cref{sec:phf}).
    \item We give a \emph{$k$-perfect} extension that admits a tunable false
    negative rate $\omega$, turning the structure into a genuine two-rate
    (second-order) Bernoulli set and a binary asymmetric channel
    (\cref{sec:k_perfect}).
    \item We lift the construction from set membership to key-value lookup, the
    \emph{Perfect Map Filter} (\gls{gls-pmf}) realizing the static \emph{Bernoulli
    map}~\cite{bmaps}, and apply it to fuzzy set queries for rank-ordered
    retrieval (\cref{sec:pmf,sec:apps}).
    \item We tie the construction's space to the Carter et al.~\cite{carter1978}
    lower bound it approaches as $\varepsilon \to 0$, situating it against the
    Bloom filter and modern succinct filters (\cref{sec:adt}).
\end{enumerate}

The Bernoulli set model~\cite{bsets} is characterized by two axioms:
\emph{element-wise independence} (membership outcomes for distinct elements are
independent) and \emph{conditional independence of rates} (the false positive
rate is independent of the objective set, given its cardinality). The
construction of the \gls{gls-phf} satisfies both axioms under standard
cryptographic hash assumptions.

The paper is organized as follows.
\Cref{sec:phf} derives the \gls{gls-phf} construction, proves its false
positive rate, and analyzes its space complexity relative to the Bloom filter.
\Cref{sec:typed_phf} extends the construction to typed elements via injections
into $\BitSet^*$.
\Cref{sec:k_perfect} introduces $k$-perfect hashing and the resulting false
negative rate.
\Cref{sec:pmf} generalizes to the approximate hash map.
\Cref{sec:apps} demonstrates applications to fuzzy set queries for
rank-ordered search.
\Cref{sec:impl} provides a C++ implementation.
\Cref{sec:adt} discusses the abstract data type of the Bernoulli set, including
space efficiency bounds and algebraic properties.

\end{document}
